\section{Physical Design}

This section discusses the physical design decisions, specifically the creation of indexes to optimise
the five workload operations identified in Section~2.

\subsection{Index Analysis}

\subsubsection{Operations 1, 2, 3, 4 -- Insert Operations}

Operations 1, 3 (\textit{insert customer}, \textit{insert team}) are simple object inserts with no
join or filter predicates; they do not benefit from additional indexes.
Operation 2 (\textit{insert booking}) reads three tables by their primary keys (\texttt{Team.Code},
\texttt{Customer.ID}, \texttt{Event\_Location.ID}) -- all of which are already indexed as primary keys.
Operation 4 (\textit{insert member}) reads \texttt{Team} by primary key and scans \texttt{Member} to
check team capacity; the scan uses the \texttt{Team} REF column.

\paragraph{Explain Plans -- Operations 1--3} \leavevmode \newline
\begin{figure}[H]
    \centering
    \includegraphics[width=\textwidth]{images/ExPlan1.png}
    \caption{Explain Plan -- Operation 1 (Register Customer)}
\end{figure}

\begin{figure}[H]
    \centering
    \includegraphics[width=\textwidth]{images/ExPlan2.png}
    \caption{Explain Plan -- Operation 2 (Add Booking)}
\end{figure}

\begin{figure}[H]
    \centering
    \includegraphics[width=\textwidth]{images/ExPlan3.png}
    \caption{Explain Plan -- Operation 3 (Add Team)}
\end{figure}

\paragraph{AutoTrace -- Operation 4} \leavevmode \newline
\begin{figure}[H]
    \centering
    \includegraphics[width=\textwidth]{images/Op2Index.png}
    \caption{AutoTrace -- Operation 4 (Add Member)}
\end{figure}

\subsubsection{Operation 5 -- Query: Get Team Members}

Operation 5 is the most frequent \textbf{read} operation (5 times/day). The query is:
\begin{lstlisting}
SELECT m.ID, m.Name, m.Surname, m.Phone, m.Email
  FROM Member m
 WHERE DEREF(m.Team).Code = :team_code;
\end{lstlisting}

This filters \texttt{Member} rows by the de-referenced team code. Without an index on \texttt{Member.Team},
Oracle must perform a full table scan over all 100 (or more) members.

\paragraph{Without Index} \leavevmode \newline
\begin{figure}[H]
    \centering
    \includegraphics[width=\textwidth]{images/OldManNOIndex.png}
    \caption{AutoTrace -- Operation 5 without index (full table scan on Member)}
\end{figure}

By inspecting the AutoTrace output we can confirm that Oracle performs a \textbf{FULL TABLE SCAN} on
\texttt{Member}, reading all rows regardless of the team predicate.

\paragraph{Index Creation} \leavevmode \newline
We create an index on the REF column \texttt{Team} of the \texttt{Member} table:

\begin{lstlisting}
CREATE INDEX idx_member_team ON Member(Team);
\end{lstlisting}

\paragraph{With Index} \leavevmode \newline
\begin{figure}[H]
    \centering
    \includegraphics[width=\textwidth]{images/OldManWithIndex.png}
    \caption{AutoTrace -- Operation 5 with \texttt{idx\_member\_team} (index range scan)}
\end{figure}

With the index in place, Oracle can use an \textbf{INDEX RANGE SCAN} and retrieve only the rows
belonging to the requested team, significantly reducing logical reads (from $\sim$100 to $\sim$5).

\paragraph{Explain Plan -- Operation 5} \leavevmode \newline
\begin{figure}[H]
    \centering
    \includegraphics[width=\textwidth]{images/ExPlan5.png}
    \caption{Explain Plan -- Operation 5 with index}
\end{figure}
